\section{Design 3.0}

\subsection{Knowledge schema}
All knowledge, the ground truth and the answers share one schema:
\begin{itemize}
  \item \textbf{Tasks:} needs, total time, goal and domain.
  \item \textbf{Steps:} one action each, an object, and used items each with a \emph{role}, the function the item serves in the step.
  \item \textbf{Step details:} destinations, settings, amounts, an \texttt{until} state, purpose, preconditions, conditions, warnings and checks.
  \item \textbf{Order:} \texttt{before} and \texttt{overlaps}.
  \item \textbf{Situations:} \texttt{has}, \texttt{lacks}, \texttt{wants}, \texttt{excludes}, time limit and count.
  \item \textbf{Differences:} \texttt{swap}, \texttt{drop}, \texttt{add\_step}, \texttt{change}, extra and dropped needs, \texttt{cannot}.
\end{itemize}

The ground truth of all \gtTasks{} tasks (\gtSteps{} steps) was re-authored in this schema:
\begin{itemize}
  \item compound teach lines were split into one action per step;
  \item every used item gained a role;
  \item each variant scenario gained situation facts and difference facts.
\end{itemize}

Answers look like this:
\begin{lstlisting}
task bake_potato
need potato
step 1 do heat oven
step 1 setting temperature 200 celsius
step 3 do prick potato
step 3 uses fork pierce
step 3 purpose steam_escapes
before 4 5
gap lacks whisk 2
\end{lstlisting}

\subsection{Architecture}
\begin{description}
  \item[Learned models.] All are averaged perceptrons trained offline in Python and run in Rust, with identical predictions (a parity test):
  \begin{itemize}
    \item a line-kind classifier;
    \item taggers for step lines, needs lines and time lines;
    \item a situation tagger;
    \item a request detector;
    \item symbol-form rewrites.
  \end{itemize}
  A line whose minimum tag margin is below a threshold $\theta$ goes to the teacher. $\theta$ is set on the development set so that kept lines are at least 95\% exact. At the end of phase A the parser kept \parserCoverage\% of the \parserDevLines{} development step lines at that precision.
  \item[Validator.] It is the last gate: only facts in the schema, about the task being taught, are stored. Each stored fact carries its provenance (source, session, parser, confidence). Corrections retract facts and never delete them.
  \item[Task identity.] IDF-weighted overlap between the request and what the agent knows of each task: its symbols and the names it learned from teaching. Each request word is weighted by how specific it is across the requests the agent has seen, which is learned, not listed.
  \item[Adaptation.] Two arms are compared:
  \begin{itemize}
    \item \emph{function}: an item the user has stands in for a lacked one if it plays the same role in some taught task; otherwise a similar stored case is transferred;
    \item \emph{rules}: patterns induced from stored cases with support from at least two tasks.
  \end{itemize}
  Both reuse a taught case of the same task first. Without support, the answer carries a \texttt{gap} line and asks the teacher, and the teacher's answer is stored as a new case.
\end{description}

\subsection{Knowledge sources}
\begin{itemize}
  \item \textbf{C1 (dialogue):} each training task is requested, taught with its teach lines, then met again with new wording, in its variants (taught after a failure) and at the end for retention.
  \item \textbf{C2 (mix):} C1 plus a generated dataset of \datasetTasks{} tasks written by Haiku in the same schema. The dataset retrains the parser and fills the case library.
\end{itemize}
